\section{Fixed-parameter pretraining generalisation bounds}
\label{sec:Fixed-parameter pretraining generalisation}

This section studies pretraining generalisation bounds for a fixed model parameter value \(\Theta\), the aim is to quantify how accurately the transition kernel \(P_\Theta\) approximates a reference transition kernel \(P^*\) when these observations may be dependent.

The dependence assumptions in this section concern the observations used during pretraining. They are logically separate from the recursive-inference chain studied in Section~\ref{sec:General state-space Markov chain framework for functional autoregressive transformers}; see Remark~\ref{rem:Inference ergodicity separate from pretraining}. In particular, the pretraining observations need not be generated by \(P_\Theta\), and the assumptions used to establish geometric ergodicity of the learned inference chain do not enter the generalisation theorem.

The analysis proceeds in four stages. We first define a transition-kernel loss and show that, because the shift part of the autoregressive transition is deterministic, this loss reduces exactly to a Wasserstein discrepancy between the reference and learned conditional output laws. We then derive a generalisation bound and a sample-complexity bound using Marton couplings. Next, we develop a complementary Wasserstein metric approach based on transportation inequalities and Wasserstein dependence. Finally, we compare the sampling regimes under which the two concentration frameworks are most naturally applied.

\subsection{Reference transition kernel, loss, and risk}
\label{subsec:Reference transition kernel loss and risk}

The statistical object to be learned is the one-step transition mechanism of the context process. We therefore begin by defining a reference context-transition kernel \(P^*\), constructed from the reference conditional law \(Q^*\) in exactly the same way that the model transition kernel \(P_\Theta\) is constructed from \(Q_\Theta\).

Let \(\mathsf S=E\times H\), where \(E=H^n\), and let \(S_i=(F_i,Y_i)\in\mathsf S\) denote the \(i\)-th pretraining observation. Since \(E\) and \(H\) are separable Hilbert spaces, \(\mathsf S\) is Polish. We assume that there is a reference conditional probability kernel \(Q^*:E\times\mathcal B(H)\to[0,1]\) such that \(Q^*(F_i,\cdot)\in\mathcal P_2(H)\) for every \(F_i\in E\), and such that, for every \(A\in\mathcal B(H)\), \(Q^*(F_i,A)=\mathbb P(Y_i\in A\mid F_i)\) almost surely. The kernel \(Q^*\) represents the reference conditional law of the next function given the context \(F_i\) and need not coincide with the model conditional law \(Q_\Theta\).

\begin{definition}[Reference context-transition kernel]
\label{def:Reference context transition kernel}
Let \(F=(f_0,\ldots,f_{n-1})\in E\). The reference context-transition kernel \(P^*:E\times\mathcal B(E)\to[0,1]\) induced by \(Q^*\) is defined by
\[
P^*(F,A)
=
\int_H
\mathbf 1_A(f_1,\ldots,f_{n-1},g)\,
Q^*(F,\mathrm dg),
\qquad
A\in\mathcal B(E).
\]
\end{definition}

Thus \(P^*(F,\cdot)\) and \(P_\Theta(F,\cdot)\) have exactly the same deterministic shift in their first \(n-1\) coordinates. They differ only through the conditional law of the newly appended function. The definition of \(P^*\) does not assert that the complete pretraining sequence \((S_i)_{i=1}^N\) is a Markov chain, the dependence between the observations \((S_i)_{i=1}^N\) is treated separately through the assumptions introduced later.
Next, to define the loss, we use the weighted context norm from Definition~\ref{def:Weighted context norm},
\(
\|F\|_\beta^2
=
\sum_{j=0}^{n-1}\beta^j\|f_j\|_H^2,\;
\beta>1,
\)
and denote by \(W_{2,\beta}\) the corresponding \(2\)-Wasserstein distance on \(\mathcal P_2(E)\).

\begin{definition}[Transition loss]
\label{def:Fixed parameter transition loss}
For a model parameter value \(\Theta\), the transition loss at context \(F\in E\) is
\(
\ell_\Theta(F)
=
W_{2,\beta}^2
\bigl(
P^*(F,\cdot),
P_\Theta(F,\cdot)
\bigr).
\)
For a pretraining observation \(S=(F,Y)\in\mathsf S\), we set \(\ell_\Theta(S)=\ell_\Theta(F)\).
\end{definition}

At first sight, Definition~\ref{def:Fixed parameter transition loss} appears to require optimal transport on the product space \(E\). The autoregressive structure makes the situation much simpler. The deterministic shift is identical under \(P^*\) and \(P_\Theta\), so all transport cost is concentrated in the newly appended coordinate. The next proposition makes this reduction exact.

\begin{proposition}[Reduction of the context-transition loss]
\label{prop:Reduction of context transition loss}
For every \(F\in E\),
\[
W_{2,\beta}^2
\bigl(
P^*(F,\cdot),
P_\Theta(F,\cdot)
\bigr)
=
\beta^{n-1}
W_2^2
\bigl(
Q^*(F,\cdot),
Q_\Theta(F,\cdot)
\bigr).
\]
Consequently,
\(
\ell_\Theta(F)
=
\beta^{n-1}
W_2^2
\bigl(
Q^*(F,\cdot),
Q_\Theta(F,\cdot)
\bigr).
\)
\end{proposition}

\begin{proof}
Fix \(F=(f_0,\ldots,f_{n-1})\in E\), and define \(T_F:H\to E\) by
\(
T_F(g)=(f_1,\ldots,f_{n-1},g).
\)
By construction,
\(
P^*(F,\cdot)=(T_F)_\#Q^*(F,\cdot)\) and 
\(
P_\Theta(F,\cdot)=(T_F)_\#Q_\Theta(F,\cdot).
\)

Let \(\pi\) be any coupling of \(Q^*(F,\cdot)\) and \(Q_\Theta(F,\cdot)\). Pushing \(\pi\) forward through \((T_F,T_F)\) gives a coupling of \(P^*(F,\cdot)\) and \(P_\Theta(F,\cdot)\). For \(g,h\in H\),
\(
\|T_F(g)-T_F(h)\|_\beta^2
=
\beta^{n-1}\|g-h\|_H^2,
\)
because the first \(n-1\) coordinates of the two contexts are identical. Hence
\[
W_{2,\beta}^2
\bigl(
P^*(F,\cdot),
P_\Theta(F,\cdot)
\bigr)
\leqslant
\beta^{n-1}
W_2^2
\bigl(
Q^*(F,\cdot),
Q_\Theta(F,\cdot)
\bigr).
\]
For the reverse inequality, let \(\widetilde\pi\) be any coupling of \(P^*(F,\cdot)\) and \(P_\Theta(F,\cdot)\). Both marginals are supported on contexts whose first \(n-1\) coordinates are \((f_1,\ldots,f_{n-1})\). Projecting \(\widetilde\pi\) onto the final coordinates therefore gives a coupling \(\pi\) of \(Q^*(F,\cdot)\) and \(Q_\Theta(F,\cdot)\), and
\[
\int_{E\times E}\|F'-G'\|_\beta^2\,\widetilde\pi(\mathrm dF',\mathrm dG')
=
\beta^{n-1}
\int_{H\times H}\|g-h\|_H^2\,\pi(\mathrm dg,\mathrm dh).
\]
Taking the infimum over all couplings gives the reverse inequality and hence equality.
\end{proof}

The proposition shows that learning the context transition kernel is equivalent, for this loss, to learning the conditional distribution of the newly appended function. The shift component of the transition is known exactly and contributes no approximation error.
Next, we recall the definitions of true and empirical risks.

\begin{definition}[Risk]
\label{def:Empirical and population transition risks}
For a model parameter \(\Theta\) and pretraining observations \(S_1,\ldots,S_N\), define the empirical risk by
\(
\widehat{\mathcal R}_N(\Theta)
=
N^{-1}\sum_{i=1}^N\ell_\Theta(S_i).
\) and the true risk by
\(
\mathcal R_N(\Theta)
=
\mathbb E
\bigl[
\widehat{\mathcal R}_N(\Theta)
\bigr].
\) The generalisation problem consists of bounding \(\ \mathcal R_N(\Theta) - \widehat{\mathcal R}_N(\Theta) \).
\end{definition}

To obtain a concentration inequality from these definitions, the loss must first be uniformly controlled. The model already has a uniform conditional second-moment bound from the inference analysis. We impose the analogous condition on the reference conditional law.

\begin{assumption}[Uniform second moment of the reference conditional law]
\label{ass:Uniform second moment reference law}
There exists \(M_2^*<\infty\) such that
\[
\sup_{F\in E}
\int_H
\|y\|_H^2
Q^*(F,\mathrm dy)
\leqslant
M_2^*.
\]
\end{assumption}

This assumption requires a uniform bound on the conditional energy of the next function. In applications, the assumption is most natural when the admissible training contexts are generated from a bounded-energy family of initial conditions, or when the reference dynamics themselves possess a uniform moment bound. A simple deterministic example is provided by the heat equation.

\begin{example}[Heat-equation data with bounded initial energy]
\label{ex:Heat equation uniform second moment}
Let \(S_t:H\to H\) denote the heat semigroup on \(H=L^2([0,1])\), and suppose that the data are generated from initial conditions satisfying \(\|u_0\|_H\leqslant R_*\) almost surely for some \(R_*<\infty\). Assume that each context \(F=(f_0,\ldots,f_{n-1})\) consists of consecutive states of the heat equation, and consider the heat equation reference conditional law \(Q^*(F,\cdot)=\delta_{S_{\Delta t}f_{n-1}}\), where \(\Delta t>0\) is the sampling interval. Then, Assumption~\ref{ass:Uniform second moment reference law} holds with \(M_2^*=R_*^2\).
\end{example}

\begin{proof}
Let \(\Pi_{R_*}^H:H\to H\) denote the metric projection onto the closed ball \(\{f\in H:\|f\|_H\leqslant R_*\}\). Since Assumption~\ref{ass:Uniform second moment reference law} is formulated on the full context space \(E\), extend the reference kernel away from the heat-equation contexts by defining \(Q^*(F,\cdot)=\delta_{S_{\Delta t}\Pi_{R_*}^H(f_{n-1})}\) for arbitrary \(F=(f_0,\ldots,f_{n-1})\in E\). For a context arising from the stated heat-equation data, the \(L^2\)-contractivity of the heat semigroup gives \(\|f_{n-1}\|_H\leqslant\|u_0\|_H\leqslant R_*\), and hence \(\Pi_{R_*}^H(f_{n-1})=f_{n-1}\). Thus this extension agrees with the reference conditional law specified in the statement of the example on every context generated by the data.

For arbitrary \(F\in E\), since \(Q^*(F,\cdot)\) is a Dirac measure,
\[\int_H\|y\|_H^2Q^*(F,\mathrm dy)=\|S_{\Delta t}\Pi_{R_*}^H(f_{n-1})\|_H^2\leqslant\|\Pi_{R_*}^H(f_{n-1})\|_H^2\leqslant R_*^2,\]
where the first inequality again uses the \(L^2\)-contractivity of the heat semigroup and the second follows from the definition of the metric projection. Taking the supremum over \(F\in E\) gives the result. Hence Assumption~\ref{ass:Uniform second moment reference law} holds with \(M_2^*=R_*^2\).
\end{proof}

Assumption \ref{ass:Uniform second moment reference law} together with the corresponding model bound, is sufficient to obtain a finite deterministic upper bound on the Wasserstein transition loss.

\begin{proposition}[Uniform boundedness of the transition loss]
\label{prop:Uniform boundedness of fixed-parameter loss}
Under Assumption~\ref{ass:Uniform second moment reference law}, the transition loss satisfies \(0\leqslant\ell_\Theta(F)\leqslant B_\Theta\) for every \(F\in E\), where
\(
B_\Theta
=
\beta^{n-1}
\left(
\sqrt{M_2^*}+\sqrt{M_2}
\right)^2,
\)
and \(M_2\) is the model second-moment constant from Lemma~\ref{lem:Uniform conditional second moment}.
\end{proposition}

\begin{proof}
Fix \(F\in E\) and let \(\delta_0\) denote the Dirac measure at the zero function. Then, we compute:
\[
\begin{aligned}
\ell_\Theta(F)
&=
\beta^{n-1}
W_2^2
\bigl(
Q^*(F,\cdot),
Q_\Theta(F,\cdot)
\bigr)
\\
&\leqslant
\beta^{n-1}
\left[
W_2
\bigl(
Q^*(F,\cdot),
\delta_0
\bigr)
+
W_2
\bigl(
\delta_0,
Q_\Theta(F,\cdot)
\bigr)
\right]^2
\\
&=
\beta^{n-1}
\left[
\left(
\int_H \|y\|_H^2\,Q^*(F,\mathrm dy)
\right)^{1/2}
+
\left(
\int_H \|y\|_H^2\,Q_\Theta(F,\mathrm dy)
\right)^{1/2}
\right]^2
\\
&\leqslant
\beta^{n-1}
\left(
\sqrt{M_2^*}+\sqrt{M_2}
\right)^2
\\
&=:
B_\Theta,
\end{aligned}
\]
The first equality follows from Proposition~\ref{prop:Reduction of context transition loss}. The first inequality follows from the triangle inequality for \(W_2\), applied with the intermediate measure \(\delta_0\). The second equality follows from Example~\ref{ex:Dirac optimal transport}. The second inequality follows from Assumption~\ref{ass:Uniform second moment reference law} for \(Q^*\) and Lemma~\ref{lem:Uniform conditional second moment} for \(Q_\Theta\). The final equality is the definition of \(B_\Theta\).
\end{proof}

Proposition \ref{prop:Uniform boundedness of fixed-parameter loss} provides the deterministic bound required by both generalisation arguments developed below. We first derive a concentration and sample-complexity result using Marton couplings. We then develop a metric alternative based on transportation inequalities and Wasserstein dependence. The sampling regimes under which each framework is appropriate are discussed jointly at the end of the section.

\subsection{Marton-coupling generalisation bounds and sample complexity}
\label{subsec:Marton coupling generalisation}

The empirical risk is an average of bounded losses, but the pretraining observations need not be independent. We therefore require a dependence notion that is weak enough to include several practically relevant sampling schemes while still supporting concentration inequalities. Marton coupling provides such a framework by measuring how a perturbation of one coordinate can propagate to later coordinates.

We follow the Hamming-coupling route used structurally by Zekri et al.~\cite{ZekriEtAl2025}. Since the observation space \(\mathsf S\) is Polish, regular conditional distributions exist and the Marton construction applies without discretising the functional observations.
The first definition describes how two versions of the future sequence are coupled after changing one observation in the past.

\begin{definition}[Marton coupling]
\label{def:Marton coupling}
Let \(\mathbf S=(S_1,\ldots,S_N)\) take values in \(\mathsf S^N\). A Marton coupling is a family of coupled pairs
\(
\left(
\mathbf S^{(s_1,\ldots,s_i,s_i')},
\mathbf S'^{(s_1,\ldots,s_i,s_i')}
\right),
\)
indexed by \(i\in[N]\) and admissible \(s_1,\ldots,s_i,s_i'\in\mathsf S\), such that the coupled sequences agree in coordinates \(1,\ldots,i-1\), have \(i\)-th coordinates fixed respectively to \(s_i\) and \(s_i'\), and their future coordinates have the corresponding regular conditional laws. If \(s_i=s_i'\), the two sequences may be chosen equal.
\end{definition}

A Marton coupling is qualitative until the probability of future disagreement is quantified. The mixing matrix records precisely this propagation of dependence: its \((i,j)\)-entry measures how strongly changing coordinate \(i\) can affect coordinate \(j\).

\begin{definition}[Marton mixing matrix]
\label{def:Marton mixing matrix}
The mixing matrix \(\Gamma\in\mathbb R^{N\times N}\) associated with a Marton coupling is the upper-triangular matrix satisfying \(\Gamma_{ii}=1\), \(\Gamma_{ij}=0\) for \(i>j\), and, for \(i<j\),
\[
\Gamma_{ij}
=
\sup_{s_1,\ldots,s_i,s_i'}
\mathbb P
\left(
S_j^{(s_1,\ldots,s_i,s_i')}
\neq
S_j'^{(s_1,\ldots,s_i,s_i')}
\right).
\]
Its Euclidean operator norm is denoted by \(\|\Gamma\|_{\mathrm{op}}\).
\end{definition}

The Marton-coupling assumption is satisfied by several standard dependence structures and is therefore non-vacuous. In the simplest case, independent observations admit a Marton coupling with mixing matrix
\(
\Gamma=I_N,
\)
so that \(\|\Gamma\|_{\mathrm{op}}=1\). Markov-chain observations also fall within the coupling-based concentration framework: Paulin derives a bounded-differences inequality for Markov chains in terms of their mixing time in \cite[Corollary~2.11]{Paulin2015}. In the present functional setting, a concrete instance is provided by independently sampled trajectories. When the initial conditions are independent and each trajectory is generated deterministically from its own initial condition, the resulting trajectory-level observations remain independent and hence have mixing matrix \(I_N\). This is verified explicitly for heat-equation trajectories in Example~\ref{ex:Heat equation Marton coupling} of Section~\ref{subsec:Training-data regimes}.
The next theorem is the concentration tool used in the pretraining analysis. We use the dependent bounded-differences theorem of Paulin \cite[Theorem~2.1]{Paulin2015}, specialised to the trivial partition.

\begin{theorem}[Dependent bounded differences under a Marton coupling]
\label{thm:Dependent bounded differences Marton}
Suppose that \(\mathbf S\) admits a Marton coupling with mixing matrix \(\Gamma\). Let \(f:\mathsf S^N\to\mathbb R\) satisfy
\(
f(s)-f(s')
\leqslant
\sum_{i=1}^N
c_i\mathbf 1_{\{s_i\neq s_i'\}}
\)
for every \(s,s'\in\mathsf S^N\), where \(c_i\geq0\). Then, for every \(t\geq0\),
\[
\mathbb P
\left(
|f(\mathbf S)-\mathbb Ef(\mathbf S)|
\geq t
\right)
\leqslant
2\exp
\left(
-\frac{2t^2}
{\|\Gamma\|_{\mathrm{op}}^2\|c\|_2^2}
\right).
\]
\end{theorem}

The theorem separates the two quantities that determine concentration. The vector \(c\) measures how sensitive the statistic itself is to each observation, while \(\Gamma\) measures how changes in one observation can propagate through the data sequence. For the empirical risk, the boundedness result above makes the first quantity explicit.
We can now apply this abstract concentration theorem to the empirical risk.

\begin{theorem}[Fixed-parameter Marton generalisation bound]
\label{thm:Fixed-parameter Marton generalisation bound}
Fix \(\Theta\), and suppose that \(\mathbf S\) admits a Marton coupling with mixing matrix \(\Gamma\). Under Assumption~\ref{ass:Uniform second moment reference law}, for every \(t\geq0\),
\[
\mathbb P
\left(
\left|
\widehat{\mathcal R}_N(\Theta)
-
\mathcal R_N(\Theta)
\right|
\geq t
\right)
\leqslant
2\exp
\left(
-\frac{2Nt^2}
{B_\Theta^2\|\Gamma\|_{\mathrm{op}}^2}
\right).
\]
Consequently, for every \(\delta\in(0,1)\), with probability at least \(1-\delta\),
\[
\mathcal R_N(\Theta)
\leqslant
\widehat{\mathcal R}_N(\Theta)
+
\frac{
B_\Theta\|\Gamma\|_{\mathrm{op}}
}{
\sqrt{2N}
}
\sqrt{
\log\left(\frac2\delta\right)
}.
\]
\end{theorem}

\begin{proof}
Define
\(
f(s_1,\ldots,s_N)
=
N^{-1}
\sum_{j=1}^N
\ell_\Theta(s_j).
\)
Let \(s=(s_1,\ldots,s_N)\) and \(s'=(s_1',\ldots,s_N')\) differ only in coordinate \(i\). Then
\(
|f(s)-f(s')|
=
N^{-1}
\left|
\ell_\Theta(s_i)-\ell_\Theta(s_i')
\right|
\leqslant{B_\Theta}/{N}
\), where the inequality follows from Proposition~\ref{prop:Uniform boundedness of fixed-parameter loss}, since both \(\ell_\Theta(s_i)\) and \(\ell_\Theta(s_i')\) lie in \([0,B_\Theta]\), and hence their difference has absolute value at most \(B_\Theta\).

Therefore the bounded-differences vector in Theorem~\ref{thm:Dependent bounded differences Marton} may be chosen as \(c=(B_\Theta/N,\ldots,B_\Theta/N)\), with norm \(\|c\|_2
=
{B_\Theta}{/\sqrt N}.
\)

Applying Theorem~\ref{thm:Dependent bounded differences Marton} now gives, for every \(t\geq0\),
\[
\begin{aligned}
\mathbb P
\left(
\left|
f(\mathbf S)-\mathbb Ef(\mathbf S)
\right|
\geq t
\right)
&\leqslant
2\exp
\left(
-\frac{2t^2}
{\|\Gamma\|_{\mathrm{op}}^2\|c\|_2^2}
\right)
\\
&=
2\exp
\left(
-\frac{2t^2}
{\|\Gamma\|_{\mathrm{op}}^2(B_\Theta^2/N)}
\right)
\\
&=
2\exp
\left(
-\frac{2Nt^2}
{B_\Theta^2\|\Gamma\|_{\mathrm{op}}^2}
\right).
\end{aligned}
\]
Noting that \(
f(\mathbf S)=\widehat{\mathcal R}_N(\Theta)\) and 
\(
\mathbb Ef(\mathbf S)=\mathcal R_N(\Theta)
\) concludes the first part.

To obtain the high-probability form, choose \(t=t_\delta\) so that the upper bound equals \(\delta\). Then
\[
\delta
=
2\exp
\left(
-\frac{2Nt_\delta^2}
{B_\Theta^2\|\Gamma\|_{\mathrm{op}}^2}
\right)
\Longleftrightarrow
\log\left(\frac2\delta\right)
=
\frac{2Nt_\delta^2}
{B_\Theta^2\|\Gamma\|_{\mathrm{op}}^2}
\Longleftrightarrow
t_\delta
=
\frac{
B_\Theta\|\Gamma\|_{\mathrm{op}}
}{
\sqrt{2N}
}
\sqrt{
\log\left(\frac2\delta\right)
}.
\]

Hence, with probability at least \(1-\delta\),
\[
\left|
\widehat{\mathcal R}_N(\Theta)
-
\mathcal R_N(\Theta)
\right|
\leqslant
\frac{
B_\Theta\|\Gamma\|_{\mathrm{op}}
}{
\sqrt{2N}
}
\sqrt{
\log\left(\frac2\delta\right)
} \implies
\mathcal R_N(\Theta)
\leqslant
\widehat{\mathcal R}_N(\Theta)
+
\frac{
B_\Theta\|\Gamma\|_{\mathrm{op}}
}{
\sqrt{2N}
}
\sqrt{
\log\left(\frac2\delta\right)
},
\]
which proves the result.
\end{proof}

The theorem gives a direct functional analogue of a dependent bounded-differences generalisation bound. Its rate is \(N^{-1/2}\), while all dependence between the observations appears through \(\|\Gamma\|_{\mathrm{op}}\). In the independent case, this factor is one and the usual bounded-differences scaling is recovered.

The preceding theorem bounds the generalisation gap for a given sample size. The complementary question is how many observations are sufficient to guarantee a prescribed population transition risk. Under the idealised assumption of zero empirical transition risk, this can be obtained directly by inverting the concentration bound.
We first formalise what is meant by pretraining sample complexity.

\begin{definition}[Pretraining sample complexity]
\label{def:Pretraining sample complexity}
For a model parameter \(\Theta\), confidence level \(\delta\in(0,1)\), and target transition risk \(\varepsilon_{\mathrm{pre}}>0\), the pretraining sample complexity is the smallest number of observations sufficient to guarantee
\(
\mathcal R_N(\Theta)
\leqslant
\varepsilon_{\mathrm{pre}}
\)
with probability at least \(1-\delta\), under the stated assumptions on the empirical risk and dependence structure.
\end{definition}

\begin{proposition}[Pretraining sample complexity under zero empirical risk]
\label{prop:Pretraining sample complexity zero empirical risk}
Fix \(\Theta\), and suppose that for every candidate sample size \(N\), the observations admit a Marton coupling with mixing matrix \(\Gamma_N\) satisfying \(\|\Gamma_N\|_{\mathrm{op}}\leqslant\Gamma_\star\) for some finite constant \(\Gamma_\star\). Let \(\varepsilon_{\mathrm{pre}}>0\) and \(\delta\in(0,1)\). If
\(
\widehat{\mathcal R}_{N_{\mathrm{train}}}(\Theta)=0
\)
and
\[
N_{\mathrm{train}}
\geq
\left\lceil
\frac{
B_\Theta^2\Gamma_\star^2
}{
2\varepsilon_{\mathrm{pre}}^2
}
\log
\left(
\frac2\delta
\right)
\right\rceil,
\]
then, with probability at least \(1-\delta\),
\(
\mathcal R_{N_{\mathrm{train}}}(\Theta)
\leqslant
\varepsilon_{\mathrm{pre}}.
\)
\end{proposition}

\begin{proof}
Theorem~\ref{thm:Fixed-parameter Marton generalisation bound}, the zero-empirical-risk assumption, and \(\|\Gamma_{N_{\mathrm{train}}}\|_{\mathrm{op}}\leqslant\Gamma_\star\) give
\[
\mathcal R_{N_{\mathrm{train}}}(\Theta)
\leqslant
\frac{
B_\Theta\Gamma_\star
}{
\sqrt{2N_{\mathrm{train}}}
}
\sqrt{
\log
\left(
\frac2\delta
\right)
}
\]
with probability at least \(1-\delta\). The stated lower bound on \(N_{\mathrm{train}}\) makes the right-hand side at most \(\varepsilon_{\mathrm{pre}}\).
\end{proof}

\begin{remark}[Nonzero empirical risk]
If instead \(\widehat{\mathcal R}_{N_{\mathrm{train}}}(\Theta)\leqslant\varepsilon_{\mathrm{train}}<\varepsilon_{\mathrm{pre}}\), the same calculation applies with required generalisation gap \(\varepsilon_{\mathrm{pre}}-\varepsilon_{\mathrm{train}}\).
\end{remark}

\subsection{Wasserstein metric generalisation bounds and sample complexity}
\label{subsec:Wasserstein and metric dependence}

The Marton-coupling result of the previous subsection measures dependence through the probability that two coupled coordinates disagree exactly. This is natural for discrete observations, but exact equality can be too rigid for function-valued data. Two trajectories may remain very close in the Hilbert norm while being different at every time point.

For example, let \(T:H\to H\) be a deterministic evolution map. If \(Tf\neq Tg\), then \(d_{\mathrm{TV}}(\delta_{Tf},\delta_{Tg})=1\), even when \(\|Tf-Tg\|_H\) is arbitrarily small. By contrast, \(W_1(\delta_{Tf},\delta_{Tg})=\|Tf-Tg\|_H\). Thus a deterministic evolution can be maximally separated in total variation while remaining stable in the natural function-space metric.

The aim of this subsection is to exploit this metric stability directly. We first impose a \(T_1\) transportation inequality on the one-step conditional laws of the data process. This controls concentration of each conditional distribution. We then impose a \(W_1\)-stability condition on the data-generating transition kernel, which controls how perturbations of the current state affect the next conditional distribution. Theorem~2.5 of Djellout, Guillin, and Wu \cite{DjelloutGuillinWu2004} combines these two properties to obtain a transportation inequality for the law of the entire dependent trajectory. This path-space inequality then yields concentration for Lipschitz empirical averages. Finally, we apply that concentration result to the transition loss.

\paragraph{One-step transportation inequality.}

We begin with the one-step concentration property used by Djellout, Guillin, and Wu. Informally, a \(T_1(C)\) inequality states that moving probability mass far from a reference distribution in \(W_1\) requires a correspondingly large relative-entropy cost. The constant \(C\) therefore determines the concentration scale of the measure.

\begin{definition}[\(T_1\) transportation inequality]
\label{def:T1 transportation inequality}
Let \((\mathsf X,d)\) be a Polish metric space. A probability measure \(\mu\in\mathcal P(\mathsf X)\) satisfies \(T_1(C)\) if, for every \(\eta\in\mathcal P(\mathsf X)\),
\(
W_1^d(\eta,\mu)
\leqslant
\sqrt{2C\,H(\eta\mid\mu)},
\)
where \(H(\eta\mid\mu)\) denotes the relative entropy of \(\eta\) with respect to \(\mu\), interpreted as \(+\infty\) when \(\eta\) is not absolutely continuous with respect to \(\mu\).
\end{definition}

Smaller values of \(C\) correspond to stronger concentration. In the dependent setting below, the same finite constant is assumed for the initial law and for every one-step conditional law. This controls the amount of randomness introduced at each step. A separate Wasserstein condition controls the propagation of perturbations between consecutive states.

\paragraph{Path space transportation inequality.}

Let \((X_i)_{i=1}^N\) be a Markov chain on \((\mathsf X,d)\), with initial distribution \(\mu_1\) and transition kernel \(K\). Denote its joint path law by \(\mathbb P_N=\mathcal L(X_1,\ldots,X_N)\), and equip \(\mathsf X^N\) with the product metric \(d_{\ell_1}(x,y)=\sum_{i=1}^N d(x_i,y_i)\).

The next proposition is the bridge between the local assumptions on the data-generating kernel and concentration of the entire dependent sequence.

\begin{proposition}[Path-space transportation inequality]
\label{prop:Path space transportation inequality}
Suppose that \(\mu_1\) and every \(K(x,\cdot)\) satisfy \(T_1(C)\) for the same finite constant \(C\). Assume further that there exists a constant \(\kappa\in[0,1)\) such that 
\(
W_1^d
\bigl(
K(x,\cdot),
K(y,\cdot)
\bigr)
\leqslant
\kappa d(x,y)
\)
for every \(x,y\in\mathsf X\).

Then the joint law \(\mathbb P_N\) satisfies \(T_1(CN/(1-\kappa)^2)\) on \((\mathsf X^N,d_{\ell_1})\). Equivalently, for every probability measure \(\mathbb Q\) on \(\mathsf X^N\),
\[
W_1^{d_{\ell_1}}
\bigl(
\mathbb Q,
\mathbb P_N
\bigr)
\leqslant
\frac{1}{1-\kappa}
\sqrt{
2CN\,H(\mathbb Q\mid\mathbb P_N)
}.
\]
\end{proposition}

\begin{proof}
For \(i\geq2\), the Markov property gives \(P_i(\cdot\mid x^{i-1})=K(x_{i-1},\cdot)\), where \(P_i(\cdot\mid x^{i-1})\) denotes the conditional law of the \(i\)-th coordinate under \(\mathbb P_N\).

Condition \emph{(C1)} of \cite[Theorem~2.5]{DjelloutGuillinWu2004} is verified by choosing \(a_1=\kappa\) and \(a_j=0\) for \(j\geq2\). Indeed,
\[
\begin{aligned}
W_1^d
\left(
P_i(\cdot\mid x^{i-1}),
P_i(\cdot\mid\widetilde x^{i-1})
\right)
&=
W_1^d
\left(
K(x_{i-1},\cdot),
K(\widetilde x_{i-1},\cdot)
\right)
\\
&\leqslant
\kappa
d(x_{i-1},\widetilde x_{i-1})
\\
&=
\sum_{j=1}^{i-1}
a_j
d(x_{i-j},\widetilde x_{i-j}).
\end{aligned}
\]
The first equality follows from the Markov property. The inequality is the assumed Wasserstein stability of \(K\). The final equality follows from the choice \(a_1=\kappa\) and \(a_j=0\) for \(j\geq2\). The dependence coefficient is therefore \(r_1=\sum_{j\geq1}a_j=\kappa<1\).

Taking \(p=1\) in \cite[Theorem~2.5]{DjelloutGuillinWu2004} gives
\[
\begin{aligned}
W_1^{d_{\ell_1}}
\bigl(
\mathbb Q,
\mathbb P_N
\bigr)
&\leqslant
\frac{1}{1-r}
\sqrt{
2CN\,H(\mathbb Q\mid\mathbb P_N)
}
\\
&=
\frac{1}{1-\kappa}
\sqrt{
2CN\,H(\mathbb Q\mid\mathbb P_N)
}
\\
&=
\sqrt{
2
\left(
\frac{CN}{(1-\kappa)^2}
\right)
H(\mathbb Q\mid\mathbb P_N)
}.
\end{aligned}
\]
Comparing the final expression with Definition~\ref{def:T1 transportation inequality} proves that \(\mathbb P_N\) satisfies \(T_1(CN/(1-\kappa)^2)\).
\end{proof}

The two hypotheses play distinct roles. The \(T_1(C)\) assumption controls concentration of each conditional distribution, whereas \(\kappa\) measures sensitivity of the next-step law to the current state. The factor \((1-\kappa)^{-1}\) records the accumulation of this dependence along the trajectory. If \(\kappa\) is small, perturbations are rapidly forgotten; as \(\kappa\) approaches one, the concentration bound deteriorates.

These assumptions are satisfied by concrete classes of stochastic dynamical systems and are therefore not merely abstract regularity conditions. For example, Djellout, Guillin, and Wu study discrete-time stochastic recursions of the form
\(
X_{k+1}=f(X_k)+\sigma(X_k)W_{k+1},
\)
for which the above assumptions are satisfied; see \cite[Section~3.2, Corollary~3.5]{DjelloutGuillinWu2004}. In the present functional setting, an explicit example is given by the heat-equation context process in Example~\ref{ex:Heat equation metric dependence} of Section~\ref{subsec:Training-data regimes}, where the \(L^2\)-contractivity of the heat semigroup provides the Wasserstein dependence coefficient and bounded initial energy yields a common finite \(T_1\) constant.

\paragraph{Concentration of empirical averages.}

Once the path law satisfies a \(T_1\) inequality, the next step is to apply it to an empirical average. Averaging is important because the Lipschitz constant of the empirical statistic is of order \(N^{-1}\), which ultimately produces the familiar \(N^{-1/2}\) concentration rate.

\begin{theorem}[Metric concentration for a Wasserstein-dependent data chain]
\label{thm:Metric concentration Wasserstein data chain}
Let \((X_i)_{i=1}^N\) satisfy the assumptions of Proposition~\ref{prop:Path space transportation inequality}. Let \(\varphi:\mathsf X\to\mathbb R\) be \(L_\varphi\)-Lipschitz, and define \(\widehat r_N=N^{-1}\sum_{i=1}^N\varphi(X_i)\) and \(r_N=\mathbb E[\widehat r_N]\). Then, for every \(t>0\),
\[
\mathbb P
\left(
|\widehat r_N-r_N|
\geq t
\right)
\leqslant
2
\exp
\left(
-\frac{
Nt^2(1-\kappa)^2
}{
2CL_\varphi^2
}
\right).
\]
Consequently, for every \(\delta\in(0,1)\), with probability at least \(1-\delta\),
\[
r_N
\leqslant
\widehat r_N
+
\frac{L_\varphi}{1-\kappa}
\sqrt{
\frac{2C}{N}
\log
\left(
\frac2\delta
\right)
}.
\]
\end{theorem}

\begin{proof}
Define \(f:\mathsf X^N\to\mathbb R\) by \(f(x_1,\ldots,x_N)=N^{-1}\sum_{i=1}^N\varphi(x_i)\). For \(x,y\in\mathsf X^N\),
\[
\begin{aligned}
|f(x)-f(y)|
&=
\left|
\frac1N
\sum_{i=1}^N
\bigl(
\varphi(x_i)-\varphi(y_i)
\bigr)
\right|
\\
&\leqslant
\frac1N
\sum_{i=1}^N
|\varphi(x_i)-\varphi(y_i)|
\\
&\leqslant
\frac{L_\varphi}{N}
\sum_{i=1}^N
d(x_i,y_i)
\\
&=
\frac{L_\varphi}{N}
d_{\ell_1}(x,y).
\end{aligned}
\]
The first inequality is the triangle inequality. The second follows from the \(L_\varphi\)-Lipschitz continuity of \(\varphi\). The final equality is the definition of \(d_{\ell_1}\). Thus \(f\) is \(L_\varphi/N\)-Lipschitz on the path space.

By Proposition~\ref{prop:Path space transportation inequality}, the path law satisfies \(T_1(CN/(1-\kappa)^2)\). Applying the concentration consequence of this transportation inequality, equivalently \cite[Corollary~2.6]{DjelloutGuillinWu2004}, gives
\[
\begin{aligned}
\mathbb P
\left(
f-\mathbb Ef
\geq t
\right)
&\leqslant
\exp
\left(
-\frac{
t^2
}{
2
\left(
CN/(1-\kappa)^2
\right)
\left(
L_\varphi/N
\right)^2
}
\right)
\\
&=
\exp
\left(
-\frac{
Nt^2(1-\kappa)^2
}{
2CL_\varphi^2
}
\right).
\end{aligned}
\]
The inequality is the concentration consequence of the path-space \(T_1\) inequality. 

Applying the same argument to \(-f\) and using a union bound gives
\[
\mathbb P
\left(
|f-\mathbb Ef|
\geq t
\right)
\leqslant
2
\exp
\left(
-\frac{
Nt^2(1-\kappa)^2
}{
2CL_\varphi^2
}
\right).
\]
Since \(f(X_1,\ldots,X_N)=\widehat r_N\) and \(\mathbb Ef=r_N\), this proves the first statement.

For the high-probability form, choose \(t=t_\delta\) such that
\[
2
\exp
\left(
-\frac{
Nt_\delta^2(1-\kappa)^2
}{
2CL_\varphi^2
}
\right)
=
\delta.
\implies
t_\delta
=
\frac{L_\varphi}{1-\kappa}
\sqrt{
\frac{2C}{N}
\log
\left(
\frac2\delta
\right)
}.
\]
\end{proof}

The theorem separates the sources of the statistical constant clearly. The factor \(N^{-1/2}\) comes from averaging \(N\) observations, \(L_\varphi\) measures sensitivity of the statistic being averaged, \(C\) measures one-step concentration, and \((1-\kappa)^{-1}\) measures persistence of dependence along the data trajectory.

\paragraph{Application to the transformer loss.}

The preceding results concern a generic Lipschitz statistic of a metric-dependent data sequence. We now specialise the statistic to the transition loss of the functional autoregressive transformer. In this part of the analysis, \(X_i=F_i\), so the observations are the functional contexts themselves. The empirical function in Theorem~\ref{thm:Metric concentration Wasserstein data chain} is therefore \(\varphi(F)=\ell_\Theta(F)\).

To apply the theorem, it remains to show that the map \(F\mapsto\ell_\Theta(F)\) is Lipschitz. Proposition~\ref{prop:Reduction of context transition loss} gives \(\ell_\Theta(F)=\beta^{n-1}W_2^2(Q^*(F,\cdot),Q_\Theta(F,\cdot))\), so regularity of the transition loss follows from Wasserstein regularity of the reference and learned conditional output laws.

\begin{proposition}[Lipschitz continuity of the transition loss]
\label{prop:Lipschitz continuity squared Wasserstein loss}
Assume that there exist \(L_*\) and \(L_\Theta\) such that
\(
W_2
\bigl(
Q^*(F,\cdot),
Q^*(G,\cdot)
\bigr)
\leqslant
L_*
\|F-G\|_\beta
\)
and
\(
W_2
\bigl(
Q_\Theta(F,\cdot),
Q_\Theta(G,\cdot)
\bigr)
\leqslant
L_\Theta
\|F-G\|_\beta
\)
for every \(F,G\in E\). Under Assumption~\ref{ass:Uniform second moment reference law},
\(
|\ell_\Theta(F)-\ell_\Theta(G)|
\leqslant
L_\ell
\|F-G\|_\beta,
\)
where
\(
L_\ell
=
2
\beta^{(n-1)/2}
\sqrt{B_\Theta}
(L_*+L_\Theta).
\)
\end{proposition}

\begin{proof}
Set \(a=W_2(Q^*(F,\cdot),Q_\Theta(F,\cdot))\) and \(b=W_2(Q^*(G,\cdot),Q_\Theta(G,\cdot))\). Then
\[
\begin{aligned}
|\ell_\Theta(F)-\ell_\Theta(G)|
&=
\beta^{n-1}
|a^2-b^2|
\\
&=
\beta^{n-1}
(a+b)|a-b|
\\
&\leqslant
\beta^{n-1}
(a+b)
\left[
W_2
\bigl(
Q^*(F,\cdot),
Q^*(G,\cdot)
\bigr)
+
W_2
\bigl(
Q_\Theta(F,\cdot),
Q_\Theta(G,\cdot)
\bigr)
\right]
\\
&\leqslant
\beta^{n-1}
(a+b)
(L_*+L_\Theta)
\|F-G\|_\beta
\\
&\leqslant
2
\beta^{(n-1)/2}
\sqrt{B_\Theta}
(L_*+L_\Theta)
\|F-G\|_\beta \\
&= L_\ell\|F-G\|_\beta.
\end{aligned}
\]
The first equality follows from Proposition~\ref{prop:Reduction of context transition loss}. The second equality uses \(a^2-b^2=(a+b)(a-b)\). The first inequality follows from the triangle inequality for \(W_2\). Indeed, inserting \(Q^*(G,\cdot)\) and \(Q_\Theta(G,\cdot)\) as intermediate measures gives an upper bound for \(a-b\); interchanging \(F\) and \(G\) gives the same bound for \(b-a\), combining the two estimates yields
\(
|a-b|
\leqslant
W_2\bigl(Q^*(F,\cdot),Q^*(G,\cdot)\bigr)
+
W_2\bigl(Q_\Theta(F,\cdot),Q_\Theta(G,\cdot)\bigr)
\). The second inequality uses the assumed Wasserstein Lipschitz bounds for \(Q^*\) and \(Q_\Theta\). Finally, Proposition~\ref{prop:Uniform boundedness of fixed-parameter loss} gives \(a,b\leqslant\beta^{-(n-1)/2}\sqrt{B_\Theta}\), yielding the final inequality.
\end{proof}

The constant \(L_*\) describes regularity of the reference conditional law, whereas \(L_\Theta\) describes regularity of the learned transformer conditional law. The latter is supplied by Proposition~\ref{prop:Lipschitz continuity of the Gaussian output kernel} in Appendix~\ref{app:Lipschitz regularity}. These constants determine sensitivity of the loss to a change in context. They are conceptually distinct from \(\kappa\), which measures temporal dependence of the data-generating process.

\begin{example}[Reference-law Lipschitz continuity for the heat equation]
Let \(S_t:H\to H\) denote the homogeneous Dirichlet heat semigroup on \(H=L^2([0,1])\), and suppose the reference next-step law is deterministic, \(Q^*(F,\cdot)=\delta_{S_{\Delta t}f_{n-1}}\). Then, the reference law is Wasserstein-Lipschitz with
\(
L_*
=
e^{-\nu\pi^2\Delta t}
\beta^{-(n-1)/2}.
\)
\end{example}

\begin{proof}
We compute
\[
\begin{aligned}
W_2
\bigl(
Q^*(F,\cdot),
Q^*(G,\cdot)
\bigr)
&=
\|S_{\Delta t}f_{n-1}-S_{\Delta t}g_{n-1}\|_H
\\
&\leqslant
e^{-\nu\pi^2\Delta t}
\|f_{n-1}-g_{n-1}\|_H
\\
&\leqslant
e^{-\nu\pi^2\Delta t}
\beta^{-(n-1)/2}
\|F-G\|_\beta.
\end{aligned}
\]
The first equality follows from the Example \ref{ex:Dirac optimal transport} Wasserstein distance between Dirac measures. The first inequality is the \(L^2\)-contractivity of the heat semigroup \(
\|S_{\Delta t}f-S_{\Delta t}g\|_H
\leqslant
e^{-\nu\pi^2\Delta t}
\|f-g\|_H,
\) and the second follows from the definition of the weighted context norm.
\end{proof}

\paragraph{Metric pretraining generalisation.}

We can now combine the data-process assumptions with the Lipschitz estimate for the transition loss. The constants \(C\) and \(\kappa\) describe the pretraining data, while \(L_\ell\) describes the statistic whose empirical average is being estimated.

\begin{theorem}[Fixed-parameter metric generalisation bound]
\label{thm:Fixed-parameter metric generalisation bound}
Fix \(\Theta\), and suppose that the context observations \((F_i)_{i=1}^N\) form a Markov chain on \((E,\|\cdot\|_\beta)\) satisfying the assumptions of Theorem~\ref{thm:Metric concentration Wasserstein data chain} with constants \(C>0\) and \(\kappa\in[0,1)\). Suppose also that Proposition~\ref{prop:Lipschitz continuity squared Wasserstein loss} holds with Lipschitz constant \(L_\ell<\infty\). Then, for every \(t\geq0\),
\[
\mathbb P
\left(
\left|
\widehat{\mathcal R}_N(\Theta)
-
\mathcal R_N(\Theta)
\right|
\geq t
\right)
\leqslant
2
\exp
\left(
-\frac{
Nt^2(1-\kappa)^2
}{
2CL_\ell^2
}
\right).
\]
Consequently, for every \(\delta\in(0,1)\), with probability at least \(1-\delta\),
\[
\mathcal R_N(\Theta)
\leqslant
\widehat{\mathcal R}_N(\Theta)
+
\frac{L_\ell}{1-\kappa}
\sqrt{
\frac{2C}{N}
\log\left(\frac2\delta\right)
}.
\]
\end{theorem}

\begin{proof}
Apply Theorem~\ref{thm:Metric concentration Wasserstein data chain} with \(\mathsf X=E\), \(d(F,G)=\|F-G\|_\beta\), and \(\varphi=\ell_\Theta\). Proposition~\ref{prop:Lipschitz continuity squared Wasserstein loss} gives \(L_\varphi=L_\ell\). Moreover, by Definition~\ref{def:Empirical and population transition risks}, \(\widehat r_N=\widehat{\mathcal R}_N(\Theta)\) and \(r_N=\mathcal R_N(\Theta)\). Substituting these identities into Theorem~\ref{thm:Metric concentration Wasserstein data chain} gives the stated two-sided concentration inequality.

For the high-probability form, choose \(t=t_\delta\) so that the right-hand side of the concentration inequality is equal to \(\delta\). Then
\[
\delta
=
2
\exp
\left(
-\frac{
Nt_\delta^2(1-\kappa)^2
}{
2CL_\ell^2
}
\right)
\Longleftrightarrow
\log\left(\frac2\delta\right)
=
\frac{
Nt_\delta^2(1-\kappa)^2
}{
2CL_\ell^2
}
\Longleftrightarrow
t_\delta
=
\frac{L_\ell}{1-\kappa}
\sqrt{
\frac{2C}{N}
\log\left(\frac2\delta\right)
}.
\]
Hence, with probability at least \(1-\delta\), we have \(|\widehat{\mathcal R}_N(\Theta)-\mathcal R_N(\Theta)|\leqslant t_\delta\). In particular, \(\mathcal R_N(\Theta)\leqslant\widehat{\mathcal R}_N(\Theta)+t_\delta\), which gives the claimed high-probability generalisation bound.
\end{proof}

The theorem has the same \(N^{-1/2}\) scaling as the Marton result, but the dependence structure is encoded differently. In the Marton theorem it appears through \(\|\Gamma\|_{\mathrm{op}}\), whereas here it appears through the one-step concentration constant \(C\) and the metric-dependence coefficient \(\kappa\).

\begin{proposition}[Metric pretraining sample complexity under zero empirical risk]
\label{prop:Metric pretraining sample complexity zero empirical risk}
Fix \(\Theta\), and suppose that the context observations satisfy the assumptions of Theorem~\ref{thm:Fixed-parameter metric generalisation bound} with constants \(C>0\), \(\kappa\in[0,1)\), and \(L_\ell<\infty\). Let \(\varepsilon_{\mathrm{pre}}>0\) and \(\delta\in(0,1)\). Assume that \(\widehat{\mathcal R}_{N_{\mathrm{train}}}(\Theta)=0\). If
\[
N_{\mathrm{train}}
\geq
\left\lceil
\frac{
2CL_\ell^2
}{
(1-\kappa)^2
\varepsilon_{\mathrm{pre}}^2
}
\log\left(\frac2\delta\right)
\right\rceil,
\]
then, with probability at least \(1-\delta\), \(\mathcal R_{N_{\mathrm{train}}}(\Theta)\leqslant\varepsilon_{\mathrm{pre}}\).
\end{proposition}

\begin{proof}
Theorem~\ref{thm:Fixed-parameter metric generalisation bound} and the zero-empirical-risk assumption give
\[
\mathcal R_{N_{\mathrm{train}}}(\Theta)
\leqslant
\frac{L_\ell}{1-\kappa}
\sqrt{
\frac{2C}{N_{\mathrm{train}}}
\log\left(\frac2\delta\right)
}
\]
with probability at least \(1-\delta\). The stated lower bound on \(N_{\mathrm{train}}\) makes the right-hand side at most \(\varepsilon_{\mathrm{pre}}\).
\end{proof}

The Marton and metric results should be interpreted as complementary concentration principles. The Marton theorem is particularly natural when the sampling design consists of independent blocks or when exact coupling probabilities are informative. The metric theorem is better adapted to a temporally dependent functional process for which nearby states remain close under the dynamics. Which theorem is more useful therefore depends on the sampling regime and the geometry of the data.

\subsection{Training-data regimes and applicability}
\label{subsec:Training-data regimes}

The Marton and metric generalisation frameworks impose different assumptions on the pretraining observations. The Marton result measures dependence through coordinatewise coupling probabilities, whereas the metric result uses the geometry of the functional state space. The appropriate theorem therefore depends both on the dynamics and on how the training observations are constructed. This distinction is particularly important for functional time series, where observations along one trajectory may be strongly dependent even when different trajectories are sampled independently \cite{HormannKokoszka2010,AueEtAl2015}.

\paragraph{Independent trajectories.}

For PDE and operator-learning datasets, a natural statistical unit is an independently generated trajectory. Suppose that \(M\) initial conditions are sampled independently and that one trajectory is generated from each. Any measurable trajectory-level observation constructed separately from these trajectories is then independent across \(m\in[M]\). Consequently, the corresponding Marton mixing matrix is
\(
\Gamma=I_M,
\)
\(
\|\Gamma\|_{\mathrm{op}}=1.
\)
The effective sample size in Theorem~\ref{thm:Fixed-parameter Marton generalisation bound} is therefore the number of independent trajectories or blocks, rather than the total number of time steps stored in the dataset.
This sampling strategy does not require temporal mixing within individual trajectories and therefore remains applicable to deterministic, periodic, or conservative dynamics.

\begin{example}[Independent heat-equation trajectories admit a Marton coupling]
\label{ex:Heat equation Marton coupling}
Let \(S_t:H\to H\) denote the heat semigroup on \(H=L^2([0,1])\). Suppose that \(M\) initial conditions \(u_0^{(1)},\ldots,u_0^{(M)}\) are sampled independently from a probability distribution \(\mu_0\) on \(H\). For each \(m\in[M]\), generate the deterministic trajectory \(u_t^{(m)}=S_tu_0^{(m)}\).
Fix a sampling time \(t_0\geq0\) and a time step \(\Delta t>0\). Define the \(m\)-th context-response observation by \(S_m=(F_m,Y_m)\in E\times H\), where
\(
F_m
=
\left(
u_{t_0}^{(m)},
u_{t_0+\Delta t}^{(m)},
\ldots,
u_{t_0+(n-1)\Delta t}^{(m)}
\right)
\)
and \(Y_m=u_{t_0+n\Delta t}^{(m)}\).
Then \(S_1,\ldots,S_M\) are independent. Consequently, they admit a Marton coupling with mixing matrix \(\Gamma=I_M\), and therefore \(\|\Gamma\|_{\mathrm{op}}=1\).
\end{example}

\begin{proof}
For each \(m\in[M]\), define the deterministic map \(\Phi:H\to E\times H\) by
\[
\Phi(u)
=
\left(
\left(
S_{t_0}u,
S_{t_0+\Delta t}u,
\ldots,
S_{t_0+(n-1)\Delta t}u
\right),
S_{t_0+n\Delta t}u
\right).
\]
Then \(S_m=\Phi(u_0^{(m)})\). Since the initial conditions \(u_0^{(1)},\ldots,u_0^{(M)}\) are independent and \(\Phi\) is deterministic, the observations \(S_1,\ldots,S_M\) are independent.

Fix \(i\in[M]\), a past \(s_1,\ldots,s_{i-1}\), and two possible values \(s_i,s_i'\) of the \(i\)-th observation. By independence, the conditional law of the future observations does not depend on the value assigned to \(S_i\). Hence
\[
\mathcal L
\left(
S_{i+1},\ldots,S_M
\,\middle|\,
S_1=s_1,\ldots,S_i=s_i
\right)
=
\mathcal L
\left(
S_{i+1},\ldots,S_M
\,\middle|\,
S_1=s_1,\ldots,S_{i-1}=s_{i-1},S_i=s_i'
\right).
\]
The two conditional future laws can therefore be coupled identically. Under this coupling, we have \(S_j^{(s_1,\ldots,s_i,s_i')}=S_j'^{(s_1,\ldots,s_i,s_i')}\) almost surely for every \(j>i\). Hence
\(
\mathbb P
\left(
S_j^{(s_1,\ldots,s_i,s_i')}
\neq
S_j'^{(s_1,\ldots,s_i,s_i')}
\right)
=
0
\)
for every \(i<j\).

By Definition~\ref{def:Marton mixing matrix}, all off-diagonal entries of \(\Gamma\) are therefore zero, while \(\Gamma_{ii}=1\). Thus \(\Gamma=I_M\), and consequently \(\|\Gamma\|_{\mathrm{op}}=1\).
\end{proof}

\paragraph{Overlapping windows from a single trajectory.}

Heavily overlapping context windows from one trajectory should not be treated as independent observations merely because they appear as separate training examples. In a deterministic system, all such windows may be functions of the same random initial condition, and adjacent windows additionally share most of their coordinates. The number of stored windows may therefore substantially exceed the number of statistically distinct sampling units.
This issue is standard in dependent functional time-series analysis: temporal dependence must be accounted for explicitly rather than removed by representing the data as separate windows \cite{HormannKokoszka2010}.

\paragraph{Block-dependent and \(m\)-dependent observations.}

When dependence is local in time, observations may instead be grouped into blocks. For \(m\)-dependent or block-dependent sequences, sufficiently separated blocks are independent, and concentration can be formulated in terms of the resulting effective number of blocks. Marton-coupling inequalities naturally accommodate such dependence structures through the mixing matrix or an appropriate partition of the observations \cite{Paulin2015}.

\paragraph{Markov and total-variation mixing data.}

If the pretraining observations form a sufficiently mixing Markov chain, coupling probabilities can be controlled by its total-variation mixing properties. In this setting, the Marton framework remains applicable; see, for example, the Markov-chain bounded-differences inequalities of Paulin \cite{Paulin2015}.
For continuous functional dynamics, however, total variation may be too coarse. If \(T:H\to H\) is deterministic and \(Tf\neq Tg\), then
\(
d_{\mathrm{TV}}(\delta_{Tf},\delta_{Tg})=1
\)
even when \(\|Tf-Tg\|_H\) is arbitrarily small. The metric framework of Section~\ref{subsec:Wasserstein and metric dependence} is designed precisely for such situations.

\paragraph{Metric-dependent functional trajectories.}

The assumptions of the metric theorem arise naturally for dissipative systems. A contractive evolution can control the propagation of perturbations in the Hilbert norm, while a bounded or sufficiently concentrated initial law can provide the required transportation inequality. Transportation inequalities for dependent dynamical systems are developed in this form by Djellout, Guillin, and Wu \cite{DjelloutGuillinWu2004}.
The heat equation gives an explicit example: its semigroup contracts in \(L^2\), which induces a Wasserstein contraction for the corresponding context process, while a bounded initial-condition distribution provides a finite \(T_1\) constant.

\begin{example}[Metric dependence for heat-equation contexts]
\label{ex:Heat equation metric dependence}
Consider the heat equation \(\partial_tu=\nu\partial_{xx}u\) on \([0,1]\), with homogeneous Dirichlet boundary conditions and diffusivity \(\nu>0\). Let \(S_t:H\to H\) denote its semigroup on \(H=L^2([0,1])\), fix a sampling interval \(\Delta t>0\), and set \(\kappa_H=e^{-\nu\pi^2\Delta t}\in(0,1)\).

Suppose that the initial condition satisfies \(\|u_0\|_H\leqslant R_*\) almost surely. For a fixed context length \(n\), let \(F_k=(u_{k\Delta t},\ldots,u_{(k+n-1)\Delta t})\in E\). Define for \(F=(f_0,\ldots,f_{n-1})\) the deterministic context update by \(\Psi(F)=(f_1,\ldots,f_{n-1},S_{\Delta t}f_{n-1})\), and let \(K(F,\cdot)=\delta_{\Psi(F)}\).

Choose \(\beta>1\) such that \(\beta^{-1}+\kappa_H^2<1\), and define \(\kappa_E=(\beta^{-1}+\kappa_H^2)^{1/2}\). Then \(K\) satisfies the inequality
\(
W_1^{\|\cdot\|_\beta}
\bigl(
K(F,\cdot),
K(G,\cdot)
\bigr)
\leqslant
\kappa_E
\|F-G\|_\beta
\)
for every \(F,G\in E\). Moreover, the initial context law and every conditional law \(K(F,\cdot)\) satisfy \(T_1(C_E)\) for a common finite constant \(C_E\). Hence the heat-equation context process satisfies the assumptions of Proposition~\ref{prop:Path space transportation inequality}.
\end{example}

\begin{proof}
We first establish contraction of the heat semigroup. Let \(w(t)=S_tf-S_tg\). Since \(w\) satisfies the homogeneous Dirichlet heat equation,
\[
\begin{aligned}
\frac{\mathrm d}{\mathrm dt}
\|w(t)\|_H^2
&=
2\nu
\int_0^1
w(t,x)\,
\partial_{xx}w(t,x)\,
\mathrm dx
\\
&=
-2\nu
\int_0^1
|\partial_xw(t,x)|^2
\,\mathrm dx
\\
&\leqslant
-2\nu\pi^2
\|w(t)\|_H^2.
\end{aligned}
\]
The first equality follows by differentiating the squared \(L^2\)-norm and using the heat equation. The second follows from integration by parts and the homogeneous Dirichlet boundary conditions. The inequality is Poincaré's inequality. Grönwall's inequality therefore gives \(\|S_tf-S_tg\|_H\leqslant e^{-\nu\pi^2t}\|f-g\|_H\), and in particular \(\|S_{\Delta t}f-S_{\Delta t}g\|_H\leqslant\kappa_H\|f-g\|_H\).

Now let \(F=(f_0,\ldots,f_{n-1})\) and \(G=(g_0,\ldots,g_{n-1})\). Then
\[
\begin{aligned}
\|\Psi(F)-\Psi(G)\|_\beta^2
&=
\sum_{j=1}^{n-1}
\beta^{j-1}
\|f_j-g_j\|_H^2
+
\beta^{n-1}
\|S_{\Delta t}f_{n-1}-S_{\Delta t}g_{n-1}\|_H^2
\\
&\leqslant
\beta^{-1}
\sum_{j=1}^{n-1}
\beta^j
\|f_j-g_j\|_H^2
+
\beta^{n-1}
\kappa_H^2
\|f_{n-1}-g_{n-1}\|_H^2
\\
&\leqslant
\left(
\beta^{-1}+\kappa_H^2
\right)
\|F-G\|_\beta^2
\\
&=
\kappa_E^2
\|F-G\|_\beta^2.
\end{aligned}
\]
The first equality follows from the definitions of \(\Psi\) and \(\|\cdot\|_\beta\). The first inequality uses the heat-semigroup estimate. The second bounds each contribution by the complete weighted context norm. The final equality is the definition of \(\kappa_E\).

Since \(K(F,\cdot)=\delta_{\Psi(F)}\), Example \ref{ex:Dirac optimal transport} gives \(W_1^{\|\cdot\|_\beta}(K(F,\cdot),K(G,\cdot))=\|\Psi(F)-\Psi(G)\|_\beta\leqslant\kappa_E\|F-G\|_\beta\). By construction, \(\kappa_E<1\).

It remains to verify the \(T_1\) condition. Let \(\mu_F\) be the law of the initial \(F_0=(u_0,S_{\Delta t}u_0,\ldots,S_{(n-1)\Delta t}u_0)\) context. The semigroup estimate and the bound \(\|u_0\|_H\leqslant R_*\) give
\[
\|F_0\|_\beta^2
=
\sum_{j=0}^{n-1}
\beta^j
\|S_{j\Delta t}u_0\|_H^2
\leqslant
R_*^2
\sum_{j=0}^{n-1}
\beta^j
\kappa_H^{2j}
=:R_E^2.
\]
Thus \(\mu_F\) is supported on a ball of radius \(R_E\), whose diameter is at most \(2R_E\).
For any probability measure \(\eta\) with finite \(H(\eta\mid\mu_F)\),
\[
\begin{aligned}
W_1^{\|\cdot\|_\beta}
(\eta,\mu_F)
&\leqslant
2R_E
d_{\mathrm{TV}}(\eta,\mu_F)
\\
&\leqslant
2R_E
\sqrt{
\frac12
H(\eta\mid\mu_F)
}
\\
&=
\sqrt{
2R_E^2
H(\eta\mid\mu_F)
}.
\end{aligned}
\]
The first inequality uses the diameter bound on the common support and the second is Pinsker's inequality. Hence \(\mu_F\) satisfies \(T_1(R_E^2)\). Finally, every conditional law \(K(F,\cdot)\) is a Dirac measure. A Dirac measure satisfies \(T_1(C)\) for every \(C>0\): finite relative entropy with respect to \(\delta_{\Psi(F)}\) forces the comparison measure to equal \(\delta_{\Psi(F)}\), in which case the Wasserstein distance is zero. Thus all conditional laws satisfy \(T_1(R_E^2)\) as well. Therefore the assumptions of Proposition~\ref{prop:Path space transportation inequality} hold with \(C_E=R_E^2\) and dependence coefficient \(\kappa_E<1\).
\end{proof}

The two heat-equation examples illustrate that the same physical system can satisfy the two statistical frameworks for different reasons. With independently sampled trajectories, the Marton result uses independence across initial conditions. Along a single dependent trajectory, the metric result instead uses dissipativity of the dynamics.

\paragraph{Nonmixing dynamics.}

Not every functional system satisfies the metric condition with \(\kappa<1\). Periodic advection may preserve \(L^2\) distances, while undamped wave dynamics can preserve an energy norm. Such systems are therefore not natural candidates for the Wasserstein-contraction data theorem.
They can nevertheless be handled by the Marton result when training consists of independently sampled trajectories. In that case, independence comes from the sampling design rather than from temporal mixing of the dynamics.

The Marton and metric results should therefore be viewed as complementary. The Marton framework is particularly natural when the sampling design produces independent or weakly dependent blocks, whereas the metric framework is better adapted to temporally dependent functional processes whose dynamics are stable in the underlying function-space metric. The appropriate generalisation theorem is therefore determined jointly by the sampling scheme and the dependence structure of the data-generating dynamics.
